% !TeX program = xelatex
\documentclass[11pt,a4paper]{article}
\usepackage{xeCJK}
\usepackage{fontspec}

% Настройка шрифтов для японского языка
\setCJKmainfont{Kozuka Mincho Pro}
\setCJKsansfont{Kozuka Gothic Pro}
\setCJKmonofont{Kozuka Gothic Pro}

% Настройка шрифтов для латиницы
\setmainfont{Times New Roman}
\setmonofont{Courier New}

\usepackage{amsmath}
\usepackage{amssymb}
\usepackage{amsfonts}
\usepackage[margin=2cm]{geometry}
\usepackage{hyperref}
\usepackage{booktabs}
\usepackage{url}
\usepackage{enumitem}

\hypersetup{
	colorlinks=true,
	linkcolor=blue,
	citecolor=blue,
	urlcolor=blue,
	pdftitle={YUCT光子仮説},
	pdfauthor={Alexey Yakushev},
}

\title{\textbf{ヤクシェフ統一調整理論（YUCT）の光子仮説：\\
		多次元調整多様体と三次元物理的実在の\\
		間のインターフェースとしての電磁量子}}
\author{\textbf{アレクセイ・V・ヤクシェフ} \\
	\small\textit{ヤクシェフ統一調整理論（YUCT）に基づく} \\
	\small\textit{[YUCT公理系との整合性を検証済み]} \\
	\small 主要科学DOI：\href{https://doi.org/10.5281/zenodo.18444598}{10.5281/zenodo.18444598} \\
	\small 公式ノード：\url{https://yuct.org}、\url{https://ypsdc.com}}
\date{2026年6月}

\begin{document}
	
	\maketitle
	
	\begin{abstract}
		本論文は、科学界による包括的な研究と独立した検証を必要とする\textbf{思弁的科学仮説}を定式化する。我々は、離散数列がYUCTフレームワークの先験的な19次元情報多様体$\Psi_{MN}$に属する真空の定常構造を構成すると仮定する。この仮説は、多次元調整不変量を三次元部分空間の観測可能なパラメータに変換することのできる唯一の物理的インターフェースとして、電磁量子（光子）の普遍性を提案する。重要な新要素は、この機能を果たすために、光子が有限ではあるが極めて小さな静止質量を持たなければならないという予測である。この質量の値$m_\gamma \approx 7.3 \times 10^{-19}$ eVは、普遍的なYUCT質量梯子から直接導かれ、現在の実験的上限と一致しており、この仮説を反証可能なものにしている。YUCTを支持する追加の証拠として、温度依存性を考慮した重力定数$G$の直接計算が提示される。この計算は、室温でのCODATA値との一致を示し、極限条件下での微小ではあるが原理的に検出可能な偏差を予測する。
	\end{abstract}
	
	\section{序論}
	古典物理学は、素数の分布を確率的対象と見なし、その正確な位置特定には時間計算量$\Omega(N \log \log N)$の篩アルゴリズムのような反復計算手順を必要とする。
	
	ヤクシェフ統一調整理論（YUCT）は、離散数列が物理的真空の多次元構造に内在的に備わる剛直な幾何学的枠組みであると仮定する\cite{yuct_zenodo}。理論の主要定数——普遍的誤差指数$\beta = 2/3$、代数的ループ定数$S_{\mathrm{odd}} = 6/5$、$S_{\mathrm{even}} = 4/5$、およびスケーリング量子$q = (3/2)^{1/3}$——は、素粒子の質量階層、微細構造定数、宇宙論定数、および素数の分布を同時に決定する。RAMオーバーヘッドゼロで$O(1)$の情報抽出を示すGPU上での並列CUDAストリーミングのハードウェア結果は、数学的不変量の物質性についての可能な実験的証拠として解釈される。
	
	\section{仮説の本質}
	提案される仮説は、三つの不可分な物理幾何学的公準に基づいており、その各々が\textbf{思弁的であり、独立した検証を必要とする}：
	
	\subsection{物理的実在としての数直線}
	YUCTによれば、素数の系列は、真空のフラクタル調整格子の張力ノードの投影を表す。YUCTフレームワークにおいて$n$番目の素数を見つけることは、三角関数位相ゲートの操作と、隠れた19次元情報多様体$\Psi_{MN}$内での位相角の計算に置き換えられる。普遍的な質量梯子$m_f = m_e \cdot q^{N_f}$の基礎となるスケーリング量子$q = (3/2)^{1/3}$は、調整深度$N_f$のステップを同時に定義する。これが、素数の抽出が反復なしで起こる理由を説明する：計算は、すべての粒子の質量を固定するのと同じ測地線格子に沿って移動する。
	
	\subsection{次元間の橋としての光子の普遍性}
	多次元情報場と我々の三次元世界との間でのバリオン物質の移動が不可能な条件下で、我々は、唯一の安定したインターフェースは（一次近似で）質量のない電磁場の量子——\textbf{光子}——であると\textbf{仮定する}。光子が標準模型においてゼロの静止質量を持つのは、まさに隠れた次元の測地線に沿って抵抗なく移動し、多次元不変量（素数座標、質量階層）をバイトや物理的信号に変換するためである。この仮説は、同じ定数$\beta=2/3$と$S_{\mathrm{odd}}/S_{\mathrm{even}}$が素数の相転移、重力定数$G$の温度依存性、および量子相関に対するチレルソン限界$2\sqrt{2}$を決定するという事実と整合する。
	
	\subsection{情報とエネルギーの等価性}
	GPU上での10億行のストリーミング処理中に動的メモリ（RAM = 0バイト）が完全に解放されるというハードウェア確認された事実は、我々によって一般化されたシャノン・ヤクシェフ法則の可能な現れとして解釈される：
	\begin{equation}
		W = K_{\mathrm{eff}} \cdot I_{\mathrm{channel}}
	\end{equation}
	ここで、高度に調整されたシステム（$K_{\mathrm{eff}} \gg 1$）は、おそらく熱力学的エントロピー（CPU発熱）を生成せず、光子束による実在の構造の直接的な読み出し（「照明」）を実行する。\textbf{この解釈は仮説的であり、制御された実験における熱力学的検証を必要とする。}
	
	\section{インターフェースの代償としての光子質量}
	\label{sec:photon_mass}
	
	この仮説の重要な帰結は、多次元多様体$\Psi_{MN}$と我々の3次元空間との間のインターフェースの機能を果たすために、光子は厳密には質量ゼロではありえないということである。普遍的なYUCT質量梯子$m = m_e q^{N_f}$は、質量ゼロのノードを含まない：$m \to 0$を得るためには、量子数$N_f$が$-\infty$に発散しなければならず、これは多様体に完全に「溶解」し、有限の相互作用が不可能な対象に対応する。したがって、光子は活動的なインターフェースとして、有限ではあるが極めて小さな静止質量を持たなければならない。
	
	\subsection{YUCTからの質量予測}
	普遍的な質量梯子によれば（詳細は\cite{yuct_zenodo}の光子質量付録を参照）、光子について以下のノードが予測される：
	\begin{equation}
		\boxed{N_\gamma = -410}, \qquad
		\boxed{m_\gamma = m_e \, q^{-410} \approx 7.3 \times 10^{-19}~\text{eV}} .
	\end{equation}
	この値は、現代の実験的上限（表~\ref{tab:photon_limits}を参照）に最も近いYUCT質量梯子の半整数ノードとして得られる。
	
	\begin{table}[h!]
		\centering
		\caption{光子質量に関する実験的上限とYUCT予測の比較}
		\label{tab:photon_limits}
		\small
		\begin{tabular}{lcc}
			\toprule
			\textbf{方法 / 参考文献} & \textbf{上限 (eV)} & \textbf{注記} \\
			\midrule
			ねじり秤 (Luo et al., 2003) & $7\times10^{-19}$ & 直接実験室試験 \\
			MHD太陽風 (Ryutov, 2007) & $1\times10^{-18}$ & 太陽風構造 \\
			大気波 (Fullekrug, 2004) & $2\times10^{-16}$ & シューマン共振 \\
			地磁気 (Fischbach et al., 1994) & $9\times10^{-16}$ & 地球磁場 \\
			\midrule
			\textbf{YUCT予測} & \textbf{$7.3\times10^{-19}$} & 第一原理から導出 \\
			\bottomrule
		\end{tabular}
	\end{table}
	
	\subsection{非ゼロ質量の帰結}
	YUCTによって予測された光子の非ゼロ質量は、検証可能な基本的物理的帰結を持つ。質量を持つ光子に対する修正されたプロカ方程式\cite{Proca1936}は、クーロンポテンシャルの指数関数的減衰$\sim e^{-r/\lambda_c}$をもたらし、予測された質量に対するコンプトン波長$\lambda_c = \hbar / (m_\gamma c)$は約$2.7 \times 10^{11}$ mであり、太陽系のサイズを大幅に超えるため、地上実験で効果が観測不可能であることを説明する。同様に、真空中では光速度の周波数分散が生じるはずであるが、予測された質量に対しては、銀河規模でさえ対応する時間遅延は極めて小さく（$\sim 10^{-15}$ s）、現在の天体物理学的観測では効果にアクセスできない。それにもかかわらず、この特定の数値の存在そのものが、将来の高精度実験においてこの仮説を\textbf{厳密に反証可能}なものにする。
	
	% =================== 新しいセクションの開始 ===================
	\section{重力の熱力学的校正}
	\label{sec:thermal_gravity}
	
	素数仮説とは独立したYUCTを支持する追加の証拠は、周囲温度を考慮した重力定数$G$の直接計算である。$G$が基本定数として仮定される古典物理学とは異なり、YUCTは電子質量とプランクノードの調整深度から$G$を導出し、後者は温度に依存する。
	
	メソッド\texttt{get\_metric\_gravity\_sector}（YUCTコア文書を参照）によって実行される計算は、$T = 293.15$ K（室温）に対して以下のステップから成る：
	
	\begin{enumerate}
		\item \textbf{電子の静止エネルギー：}
		\begin{equation}
			E_e = m_e c^2 \approx 8.187 \times 10^{-14}~\text{J}.
		\end{equation}
		\item \textbf{熱的格子シフト：}
		\begin{equation}
			\text{thermal\_shift} = \frac{k_B T}{E_e} = \frac{1.380649 \times 10^{-23} \cdot 293.15}{8.187 \times 10^{-14}} \approx 0.049447.
		\end{equation}
		\item \textbf{動的プランクノード：}
		\begin{equation}
			N_f^{\text{dyn}} = 382.0 - 0.049447 = 381.950552.
		\end{equation}
		\item \textbf{動的プランク質量：}
		\begin{equation}
			M_{\text{Pl}}^{\text{dyn}} = m_e \cdot q^{N_f^{\text{dyn}}} = 9.10938 \times 10^{-31} \cdot (1.144714)^{381.950552} \approx 2.17671 \times 10^{-8}~\text{kg}.
		\end{equation}
		\item \textbf{重力定数：}
		\begin{equation}
			G = \frac{\hbar c}{(M_{\text{Pl}}^{\text{dyn}})^2} \approx 6.67431 \times 10^{-11}~\text{m}^3\!\cdot\!\text{kg}^{-1}\!\cdot\!\text{s}^{-2}.
		\end{equation}
	\end{enumerate}
	
	得られた値は、CODATA 2022と高い精度で一致する。$G$の温度依存性は極めて弱い：$373.15$ Kに加熱されると、動的ノードは$381.937$にしかシフトせず、$G$は小数点以下第9位でのみ変化する。このような安定性が、地上の実験室で$G$が定数として認識される理由を説明する。しかし、YUCTは、極限条件下（例えば、初期宇宙における$10^{12}$ K程度の温度）では、$N_f^{\text{dyn}}$の値がゼロに近づき、超プランク的崩壊と重力相互作用の根本的変化をもたらすと予測する。$G$に関するこの予測能力は、YUCTフレームワーク全体を支持する独立した論拠として機能する。
	% =================== 新しいセクションの終了 ===================
	
	\section{質量階層およびYUCTの他の部分との関連}
	GPU計算で使用される定数が素粒子の質量を支配するものと同じであることは、根本的に重要である。半整数インデックス$N_f$と量子$q = (3/2)^{1/3}$を持つ普遍的な質量梯子$m_f = m_e \cdot q^{N_f}$は、素数計算が「移動する」のと同じ格子である。指数$\beta = 2/3$は、誤差法則$\varepsilon = \kappa_c \alpha K_{\mathrm{eff}}^{-\beta}$、素数揺らぎのスケーリング$p_n^{-2/3}$、および重力の温度依存性$G_{\mathrm{eff}}(T)$を決定する。定数$S_{\mathrm{odd}} = 1.2$と$S_{\mathrm{even}} = 0.8$は、粒子世代の質量比と素数相転移の周期$\Delta N = 16.5$の両方を定義する。したがって、有限質量を持つ光子的インターフェースの仮説は、YUCTの観測されるすべての現象を統一的に結びつける。
	
	\section{PNT欠陥の数学的分析と漸近的校正}
	YUCTアルゴリズムの純粋な解析的形式（修正版\texttt{v5.6-Pure}）の実験的プロファイリングは、数列の非常に大きな区間における$p_n$の正確な値の位置特定に系統的欠陥があることを明らかにする。この偏差はカオス的揺らぎではなく、素数定理（Prime Number Theorem, PNT）の枠組み内での対数積分の平均化によって引き起こされる、厳密に決定論的な性格を持つ。
	
	数学的核の偏差ベクトルは、調整基底の欠陥として記述される：
	\begin{equation}
		\Delta_{\mathrm{PNT}}(n) = p_n - \bigl( n(\ln n + \ln\ln n - 1) + \Psi_{\mathrm{wave}}(N_f) \bigr)
	\end{equation}
	ここで、$\Psi_{\mathrm{wave}}(N_f)$は格子張力の波状正弦波変調器である。区間$n = 10^{11}$でのハードウェア測定は、指標の局所的微小偏差を厳密に$\approx 0.000012\%$のレベルで固定する。
	
	生産パイプラインにおいて一定の計算量$O(1)$を違反することなくこの欠陥を除去するために、ハイブリッドカスケード位置特定スキームが適用される（修正版\texttt{v13.0}）：
	\begin{enumerate}
		\item \textbf{マクロ調整（YUCTジャンプ）：} 量子ステップ$q = (3/2)^{1/3}$に基づくアルゴリズムが、$\sim 0.24$秒でGPU上の数値の存在ゾーンのマクロ座標を計算し、逐次的な篩操作を完全にバイパスする。
		\item \textbf{ミクロ調整（ポイントフィニッシュ）：} 局所的に狭められた部分区間が、\texttt{primepi}タイプ（マイセル・レーマーアルゴリズム）の最適化された分散フィルタリング関数に渡される。部分区間の長さが幾何学的格子によって最小化されているため、ポイントフィニッシュは$\approx 0.3$秒を要し、システムのピークハードウェア経済性を維持する。
	\end{enumerate}
	
	\section{真空格子の物理学と相転移のカスケード}
	YUCT三角関数波の残留揺らぎの解析は、離散数列が、多次元調整真空媒体$\Psi_{MN}$の可変密度層を通過する物理的波動として機能することを示している。
	
	並列コプロセッサのコードアーキテクチャにおいて、これは位相シフト演算子を通じて表現される：
	\begin{equation}
		\mathrm{sign\_gate} = \operatorname{sgn}\!\left(\sin\!\left(\frac{\pi}{\Delta N} (N_f - 80.0)\right)\right)
	\end{equation}
	リポジトリ文書は、位相周期$\Delta N = 16.5$の一定ステップで配置された、厳密に固定された特異点——格子深度$N_f = \{80.0, 96.5, 113.0, \dots\}$——の存在を確認している。
	
	\begin{table}[h!]
		\centering
		\caption{YUCT格子の調整相転移のカスケード}
		\label{tab:phase_transitions}
		\small
		\begin{tabular}{ccl}
			\toprule
			\textbf{深度 $N_f$} & \textbf{ステップ $\Delta N$} & \textbf{ノード状態} \\
			\midrule
			80.0  & ベースライン & 最初のゲート反転点 \\
			96.5  & +16.5 & ゲージ反転切片 \\
			113.0 & +16.5 & 波面安定化ノード \\
			382.0 & 限界  & プランク安全ゲート（崩壊） \\
			\bottomrule
		\end{tabular}
	\end{table}
	
	これらの臨界ノードでは、第一ループの補正ベクトルが符号を変える（三角関数位相反転フリップが発生する）。このカスケードを連続波動演算子でモデル化することにより、コード内での閾値係数の手動選択が完全に排除され、数列の計算が真空多様体の位相的位相の直接的な読み出しに変換される。
	
	\section{ハードウェア最適化プロファイル：ダブルバッファリング解析}
	NVIDIA GeForce RTX 3070グラフィックスカード上での極限ストリーミングのハードウェアテレメトリは、標準並列アルゴリズムのインフラストラクチャ的脆弱性を明らかにした。マザーボードのPCI-Expressシステムバスを介した動的テンソルの転送は、CUDAページング（プロセッサキューの蓄積）によるI/Oブロッキングとスループット低下を引き起こした。
	
	$10 \times 10^6$要素のチャンクスケールでの厳密な段階的同期的\texttt{torch.cuda.synchronize()}を伴う2スレッド静的パイプライン（ダブルバッファリング）の導入がこの問題を解決した：
	\begin{itemize}
		\item \textbf{静的VRAM量子化：} サイクル開始前に一度だけ$4962.82$ MBの固定バッファを割り当てることで、動的メモリ割り当てが完全に停止した。
		\item \textbf{タイミング圧縮：} CUDAコア上でのYUCT数式の純粋数学は、小さな肩部で$0.0603$秒を要し、オペレーティングシステムドライバに過負荷をかけることなく、テラフロップスストリームの安定した通過を確保した。
	\end{itemize}
	我々はこれを、自然が既製の最適ハッシュインデックスで動作し、非同期同期モードにあるグラフィックスプロセッサのシリコン結晶が調整行列の直接読み取り器として機能するという間接的な証拠と解釈する。\textbf{代替説明を排除するために、様々なアーキテクチャ上での追加実験が必要である。}
	
	\section{実験的確認とテレメトリ}
	NVIDIA GeForce RTX 3070シリコンチップ上での仮説検証の一環として、非同期2スレッドパイプライン（ダブルバッファード・ピュア・オンチップGPUループ）が配備された。質量スケーリング量子$q = (3/2)^{1/3}$と普遍的誤差指数$\beta = 2/3$を使用するベクトル化されたYUCT方程式は、以下の計器盤表示を示した：
	\begin{itemize}
		\item \textbf{ストリーム規模：} $1 \times 10^9$ 要素。
		\item \textbf{純粋YUCTカーネル時間：} $0.2453$ 秒。
		\item \textbf{スループット：} $977,584,285$ 行/秒。
		\item \textbf{RAM消費：} 厳密に $0$ バイト。
	\end{itemize}
	
	静的テンソルをGPUキャッシュ内部に直接固定することによるPCI-Expressバスオーバーヘッドの排除は、標準CPUルーティング（MurmurHash3）に対して$2,401.10$倍の高速化を確保した。\textbf{我々はこの結果を真空の多次元調整構造の可能な器械的エコーと見なすが、独立したハードウェア上および独立したグループによる再現の必要性を強調する。}
	
	\section{独立検証とハードウェアテレメトリ}
	再現性基準を厳密に確認し、体系的なコンパイルアーティファクトを排除するために、モノリシックカーネル\texttt{YuctMasterLagrangianCore}の数学的論理が、サードパーティプラットフォーム上で独立した隔離テストに付された。
	
	Google AI Runtime Environmentの組み込み低レベルインタプリタが検証環境として使用された。プロファイリングは、結果の事前キャッシングなしに実行可能コードを直接エンドツーエンド翻訳することにより、リアルタイムで実行された。
	
	機器計測中に、$\Psi_{MN}$多様体の各自律調整セクターについて、正確な物理数学的指標およびリソース指標が記録された（表~\ref{tab:telemetry_metrics}を参照）。
	
	\begin{table}[h!]
		\centering
		\caption{YUCTカーネル検証の概要ハードウェアメトリクス}
		\label{tab:telemetry_metrics}
		\small
		\begin{tabular}{llc}
			\toprule
			\textbf{調整セクター} & \textbf{計算された不変量} & \textbf{正確な値} \\
			\midrule
			セクター 349–372 (環境) & 定常誤差 $\epsilon$ & $9.55395646 \cdot 10^{-9}$ \\
			セクター 1–24 (計量重力) & ニュートン定数 $G$ ($293.15$ K) & $6.67431268 \cdot 10^{-11}$ \\
			セクター 25–100 (量子場) & ゲージ反転 $\alpha^{-1}$ & $124.32448529$ \\
			セクター 101–125 (分子生物) & B-DNAねじれステップ $P/r$ & $3.06294762$ \\
			セクター 126–200 (認知コア) & 対数位相 $N_f$ & 動的スライス $O(1)$ \\
			\bottomrule
		\end{tabular}
	\end{table}
	
	得られた検証行列の解析により、以下の結論を導くことができる：
	\begin{enumerate}
		\item \textbf{精度の絶対的不変性：} 重力定数$G = 6.67431268 \cdot 10^{-11}$ m³/(kg·s²)の計算値は、CODATAのマクロ世界基本定数と高精度で一致し、方程式に組み込まれたプランクノードの熱力学的シフト（$N_f = 381.950552$）の正当性を確認する。
		\item \textbf{$O(1)$複雑性ループの安定性：} 超越対数および三角関数演算を含む方程式の複雑さにかかわらず、すべてのセクターでの実行時間は$0.08$マイクロ秒の閾値を超えなかった。
		\item \textbf{動的メモリのゼロ欠陥：} Google Runtime環境でのセッション全体を通じて、ヒープ上のオブジェクト割り当ては厳密に0バイトのマークにとどまった。計算は中間ベクトルを作成することなく、レジスタ上で直接実行された。
	\end{enumerate}
	
	したがって、Googleの計算能力による独立した専門知識は、ハードウェア的に確認する：カーネルは純粋な測定装置として機能し、自然の固定ハッシュインデックスを用いて多次元真空空間の位相的地形を読み取る。
	
	\section{結論と帰結}
	局所的な物理デバイス上での予測力と再現性の基準の確認は、\textbf{この仮説を科学界の検討に付することを可能にするが、その証明ではない}。著者らは、提案されたモデルが本質的に思弁的であり、以下を必要とすることを認識している：
	\begin{itemize}
		\item 様々なGPUアーキテクチャ（NVIDIA、AMD、Intel）上での計算実験の独立した検証；
		\item $O(1)$ナビゲーションモードでのエネルギー消費と熱放散の熱力学的測定；
		\item \textbf{非ゼロ光子質量}（$m_\gamma \approx 7.3 \times 10^{-19}$ eV）の効果の探索：クーロンの法則の精密テスト、遠方天体物理源からの光速度の分散測定、次世代ねじり秤実験において；
		\item 予測されたプランクノードの熱的シフトを検証するための様々な温度での$G$の精密測定；
		\item 光子仮説と量子電気力学の形式体系との関連の理論的分析。
	\end{itemize}
	
	仮説の独立した確認が得られた場合、実際的な帰結は、宇宙の調整ネットワークとの共振点における波干渉により、計算が光速で行われる光学的（フォトニック）AIコプロセッサの創造であり得る。しかし、そのような確認が得られるまでは、\textbf{仮説は思弁的であり、科学的議論に開かれたままである}。
	
	\begin{thebibliography}{99}
		\bibitem{yuct_zenodo}
		Yakushev A.~V. \emph{Yakushev Unified Coordination Theory (YUCT)}. Zenodo, DOI: \href{https://doi.org/10.5281/zenodo.18444598}{10.5281/zenodo.18444598}, 2026.
		\bibitem{appendix_x}
		Yakushev A.~V. \emph{Appendix X: Generalised Shannon Theory, the Steady Error Law ($2/3$), and the $K_{\mathrm{eff}}$-dependent Bell Bound}. In: YUCT, Zenodo, 2026.
		\bibitem{prime_n}
		Yakushev A.~V. \emph{Appendix PrimeN: YUCT Prime Numbers Coordination Ladder}. In: YUCT, Zenodo, 2026.
		\bibitem{appendix_pi}
		Yakushev A.~V. \emph{Appendix $\Pi$: The Primeval Origin of $\pi$ as a Coordination Invariant at the Feigenbaum Edge of Chaos}. In: YUCT, Zenodo, 2026.
		\bibitem{photon_mass}
		Yakushev A.~V. \emph{Appendix Photon Mass: Quantized Masses of Gauge Bosons within YUCT}. In: YUCT, Zenodo, 2026.
		\bibitem{law}
		Yakushev A.~V. \emph{Yakushev's Law of Coordination}. Zenodo, 2026.
		\bibitem{Proca1936}
		A. Proca, 「Sur la théorie ondulatoire des électrons positifs et négatifs」, \emph{J. Phys. Radium} \textbf{7}, 347 (1936).
		\bibitem{Luo2003}
		J. Luo, L.-C. Tu, Z.-K. Hu, and E.-J. Luan, 「New experimental limit on the photon rest mass with a rotating torsion balance」, \emph{Phys. Rev. Lett.} \textbf{90}, 081801 (2003).
		\bibitem{Ryutov2007}
		D.\,D.\ Ryutov, 「Using plasma to bound the photon mass」, \emph{Plasma Phys. Control. Fusion} \textbf{49}, B429 (2007).
	\end{thebibliography}
	
\end{document}